\documentclass[../../../../main.tex]{subfiles}
\begin{document}

\begin{flushright}
\footnotesize\textit{【自動文字起こし・要確認】}
\end{flushright}

\noindent
[\textbf{解}]
\[
\int_1^x (x - t) f(x) \, dt = x^4 - 2x^2 + 1 \dots \text{①}
\]
\begin{align*}
(\text{左辺}) &= x f(x) \int_1^x dt - f(x) \int_1^x t \, dt \\
&= x f(x) (x - 1) - f(x) \frac{1}{2} (x^2 - 1) \\
&= \frac{1}{2} f(x) (x - 1)^2
\end{align*}
\[
(\text{右辺}) = (x + 1)^2 (x - 1)^2
\]
だから、①が $1 < x$ の全ての $x$ で成立する時
\[
f(x) = 2(x + 1)^2
\]
であるから、$f(x) = 2(x + 1)^2$ \hfill $\mathbin{/\mkern-5mu/}$

\bigskip

\begin{center}
\begin{tabular}{|l|}
\hline
問題文の $f(x)$ が $f(t)$ の可能性が有り、その場合は微分をくり返せばよい。\\
$\Rightarrow$ 確認したところ、$f(t)$ だったので以下に解答をのせておく。\\
\hline
\end{tabular}
\end{center}

\bigskip

\noindent
[\textbf{解2}]
\[
\int_1^x (x - t) f(t) \, dt = x^4 - 2x^2 + 1
\]
\[
\iff \int_1^x f(t) \, dt + x f(x) - x f(x) = 4x^3 - 4x \quad (x = 1 \text{で与式が成立})
\]
両辺 $x$ で微分して
\[
f(x) = 12x^2 - 4 \quad \mathbin{/\mkern-5mu/}
\]

\end{document}
